\section{Divisor channels, density one, and two distinct arithmetic maps}
\label{sec:esv3-audit}

This dated continuation (4 September 2026) checks the supplied Pro-mode
\emph{Divisor Channels, Prime Saturation, and B\'ezout Openings}, version 3,
and extends its final arithmetic comparison. The original packet is retained
unchanged under \path{supporting_materials/workbench/evidence/koide_es_v3_audit_20260904/source/}.
The claims below concern explicit reciprocal equations, integral matrices,
and quadratic surfaces. They do not require a complex structure on a sphere.

\subsection{The complete oriented channel map}

Fix integers $n\equiv1\pmod4$, $R\equiv3\pmod4$, $0<R<n$,
$\gcd(n,R)=1$, and put $a=(n+R)/4$, $N=na$. Then
$\gcd(n,a)=\gcd(R,a)=1$. For a positive divisor $\delta\mid n^2$ define
\[
 \mathcal C_\delta(n,R)=\{(u,v)\in\mathbb Z_{>0}^2:
 uv=a^2,\ R\mid4u+\delta\}.
\]
\begin{proposition}\label{prop:esv3-channel}
The following is a bijection onto all ordered positive solutions of
$4/n=1/a+1/y+1/z$:
\[
 (\delta,u,v)\longmapsto
 \left(y,z\right)=
 \left(\frac{N+(n^2/\delta)u}{R},\frac{N+\delta v}{R}\right),
 \qquad (u,v)\in\mathcal C_\delta(n,R).
\]
Its inverse, retaining this orientation, is
\[
 d=Ry-N,\quad e=Rz-N,\quad k=\gcd(d,n^2),\quad
 \delta=n^2/k,\quad u=d/k,\quad v=e/\delta.
\]
Exchange of $y,z$ is the fixed-point-free involution
$(\delta,u,v)\mapsto(n^2/\delta,v,u)$.
\end{proposition}
\begin{proof}
The reciprocal equation is $Ryz=N(y+z)$, hence $de=N^2$.
For any positive solution $Ry=N+Ny/z>N$, so $d>0$; similarly $e>0$.
Because $n$ and $a$ are coprime, every prime dividing $n$ occurs in $de$
with precisely its exponent in $n^2$. Thus $k$ extracts the entire
$n$-primary part of $d$, $\delta$ the entire $n$-primary part of $e$,
and $u,v$ are positive integers with $uv=a^2$.
Reducing $N+e$ modulo $R$ gives
$0\equiv na+\delta v\equiv4a^2+\delta v=v(4u+\delta)$.
Since $v$ is a unit modulo $R$, this proves the channel condition.
Conversely that condition makes $N+e$ divisible by $R$; also
$\delta(N+d)=n(\delta a+nu)\equiv na(\delta+4u)\equiv0\pmod R$.
The unit $\delta$ can be cancelled modulo $R$, proving integrality of $y$.
The identity $de=N^2$ proves the reciprocal equation. These operations are
inverse. A fixed point of the involution would have $\delta=n$, $u=v=a$,
so $R\mid4a+n=2n+R$, impossible since $R\ge3$ is odd and coprime to $n$.
\end{proof}

The package therefore corrects the earlier orientation reversal: the channel
order is $n^2/\gcd(Ry-na,n^2)$, not $\gcd(Ry-na,n^2)$.
For $n=9$, $R=7$, $a=4$, $(\delta,u,v)=(3,1,16)$, the ordered
solution is $(y,z)=(9,12)$, and $\gcd(27,81)=27$.

\subsection{The integral presentation, with its exact realization boundary}

Put $T=(4u+\delta)/R$, $j=(R-3)/4$, and
$\sigma=1$ for $\delta\equiv1\pmod4$, $\sigma=-1$ otherwise. Define
\[
 \ell_0=(T+\sigma)/4,\quad \ell_1=u-jT=(3T-\delta)/4,
 \quad \ell_2=\sigma,\quad
 \mathsf R_\delta=\begin{pmatrix}3&-\ell_1\\4&-T\end{pmatrix}.
\]
All coordinates are integral because $T\equiv-\delta\pmod4$.
Direct multiplication and subtraction give
\[
 \begin{pmatrix}1&-j\\0&1\end{pmatrix}
 \begin{pmatrix}R&-u\\4&-T\end{pmatrix}=\mathsf R_\delta,
 \qquad12\ell_0-4\ell_1-3\ell_2=\delta.
\]
The entries $3,4$ are coprime, so
$\operatorname{SNF}(\mathsf R_\delta)=\operatorname{diag}(1,\delta)$;
in particular its cokernel is $\mathbb Z/\delta\mathbb Z$.
Reducing $4\ell_1=3T-\delta$ modulo three gives
$\ell_1\equiv-\delta\pmod3$. Thus the source gluing admissibility
$3\nmid\ell_1$ fails precisely when $3\mid\delta$.
For the example above $(\ell_0,\ell_1,\ell_2)=(0,0,-1)$.
The matrix and its cokernel remain valid, but assigning the source's global
manifold homology to this nonadmissible triple is not established by this
calculation. This is an explicit admissibility test, not a failure of the
reciprocal solution or the integral matrix morphism.

\subsection{An exact density-one statement and what it leaves out}

\begin{proposition}\label{prop:esv3-diagonal}
Let $p\equiv1\pmod4$ be prime and $m=(p+1)/2$. The exterior channel
states on $u=a=v$ are indexed exactly by
\[
 \{R:R\mid m,\ R\equiv3\pmod4\}.
\]
For each such $R$, set $x=(p+R)/4$ and $k=(p+1)/R$. Then
\[
 \frac4p=\frac1x+\frac1{xk}+\frac1{pxk},\qquad
 \nu(p)=\frac12\left(\tau(m)-\sum_{d\mid m}\chi_4(d)\right),
\]
where $\chi_4(d)=1$ for $d\equiv1\pmod4$ and $-1$ for $d\equiv3\pmod4$.
The count is zero exactly when every prime divisor of $m$ is $1\pmod4$.
\end{proposition}
\begin{proof}
For $\delta=1$ and $u=a$, the condition is $R\mid4a+1=p+R+1$,
equivalently $R\mid p+1$. Since $R$ is odd this is $R\mid m$.
Every such $R$ obeys $3\le R\le(p+1)/2<p$, and $k\ge2$.
Adding the final two fractions gives $(p+1)/(pxk)=R/(px)$,
and adding $1/x$ gives $(p+R)/(px)=4/p$.
Distinct $R$ give distinct $x$. Summing the indicator
$(1-\chi_4(d))/2$ over divisors of $m$ proves the count.
A divisor of residue three exists exactly when a prime factor of residue
three exists: the converse follows because products of residue-one primes
and their powers have residue one.
\end{proof}

\begin{proposition}\label{prop:esv3-density}
For fixed $4\mid M$, $\gcd(c,M)=1$, $c\equiv1\pmod4$, the preceding
diagonal states cover relative density one of primes $p\equiv c\pmod M$.
For a finite set $\mathcal Q$ of primes $q\equiv3\pmod4$, $q\nmid M$,
the subfamily with some $q\mid p+1$ has exact relative density
\[
 1-\prod_{q\in\mathcal Q}\frac{q-2}{q-1}.
\]
\end{proposition}
\begin{proof}
The Chinese remainder theorem gives $\prod(q-1)$ reduced residue classes
modulo $M\prod q$ above the specified class modulo $M$. Avoiding
$p\equiv-1\pmod q$ at every $q$ leaves exactly $\prod(q-2)$ of them.
The prime number theorem for each fixed arithmetic progression gives the
displayed ratio. For completeness, its modulus-four case and partial
summation give
\[
 \sum_{\substack{q\le Y\\q\equiv3\ (4)}}\frac1q
 =\frac{\pi(Y;4,3)}Y+\int_2^Y\frac{\pi(t;4,3)}{t^2}\,dt
 \sim\tfrac12\log\log Y.
\]
Removing finitely many primes dividing $M$ does not change divergence.
Since $\log(1-z)\le-z$ for $0<z<1$, the finite products tend to zero.
The full family's lower relative density is at least each of these finite
union densities, hence is at least one and therefore equals one.
No interchange of two unproved asymptotic limits is needed.
\end{proof}

These analytic tools are classical. Elsholtz and Tao,
\emph{Counting the number of solutions to the Erd\H{o}s--Straus equation on
unit fractions}, J. Aust. Math. Soc. 94 (2013), 50--105,
\href{https://arxiv.org/abs/1107.1010v6}{arXiv:1107.1010v6}, Appendix A,
Lemma A.3, give a stronger shifted-prime additive-function statement,
attributed there to Barban via Elliott. Choosing its prime weight to be
the indicator of residue three gives the relevant abundance of factors.
The finite-CRT proof above records the precise chosen ray and progression;
it is not a claim that almost-all solvability is new.

For $p=73$, $(p+1)/2=37$ has no residue-three prime factor. This prime
therefore misses the diagonal ray, although
\[
 \frac4{73}=\frac1{20}+\frac1{220}+\frac1{4015}.
\]
Here $R=7$, $u=5$, $v=80$, and the primitive factor ray is $[1:4]$.
Thus a prime missed by the selected ray need not be missed by the full
selector. Density one does not settle the universal Erd\H{o}s--Straus claim.

\subsection{The B\'ezout opening, with the positivity hypothesis retained}

For coprime positive integers $r,t$, choose $x,y\in\mathbb Z$ with
$ry-tx=1$, set $v=(r,t)^{\mathsf T}$, $w=(x,y)^{\mathsf T}$, and let
\[
 B=(v\ w)\in\operatorname{SL}_2(\mathbb Z),\qquad
 A=\delta vv^{\mathsf T}-ww^{\mathsf T}
 =B\operatorname{diag}(\delta,-1)B^{\mathsf T}.
\]
Then $\det A=-\delta$ for every real $\delta$. All companions are
$w+kv$ for $k\in\mathbb Z$: subtraction of two B\'ezout equations gives
$r\Delta y=t\Delta x$, so coprimality gives $(\Delta x,\Delta y)=k(r,t)$.
Choosing $|x|\le r/2$ retains an actual companion rather than quotienting
away its dependence. For a positive integer channel order,
\[
 A_{11}=\delta r^2-x^2\ge(\delta-1/4)r^2>0.
\]
The broader wording ``every $\delta>0$'' does not imply that inequality.
At $(r,t,\delta)=(2,1,1/8)$ all companions have $x$ odd, and
$A_{11}=1/2-x^2\le-1/2$. For a specified real order and companion the
exact necessary and sufficient condition is $\delta>(x/r)^2$.
This correction leaves every positive integral divisor channel intact.

With $A=\left(\begin{smallmatrix}a&b\\b&c\end{smallmatrix}\right)$ and
$a>0$, the periods $(\tau,\mu,\beta)=(ia,ib/\sqrt6,ic)$ give
\[
 P_A=\begin{pmatrix}
 0&0&1&0\\ \sqrt6 b&a&0&0\\0&0&0&1\\c&b/\sqrt6&0&0
 \end{pmatrix},\qquad
 \det P_A=ac-b^2=-\delta.
\]
This invertible real matrix sends $\mathbb Z^4$ to a discrete cocompact
lattice in $\mathbb C^2$ and therefore defines a complex two-torus.
The determinant sign and the factor $\sqrt6$ are retained.
It does not determine a point on the particular global equivariant period
curve of the proposed threefold. That incidence problem remains a distinct
map problem, not an obstruction to the torus just constructed.

There is also an explicit integral equivalence of presentations, not merely
an equality of determinants. Set
\[
 P_0=\begin{pmatrix}-1&1\\-4&3\end{pmatrix},\quad
 Q_0=\begin{pmatrix}1&T-\ell_1\\0&1\end{pmatrix},\quad
 J=\begin{pmatrix}0&1\\1&0\end{pmatrix}.
\]
Multiplication gives $P_0\mathsf R_\delta Q_0=\operatorname{diag}(1,-\delta)$
and $(-JB^{-1})A(B^{-\mathsf T}J)=\operatorname{diag}(1,-\delta)$.
Thus, for $L=P_0^{-1}(-JB^{-1})$, $V=B^{-\mathsf T}JQ_0^{-1}$,
\[
 \mathsf R_\delta=LAV.
\]
As maps $\mathbb Z^2\to\mathbb Z^2$, the commuting square has source
isomorphism $V^{-1}$ and target isomorphism $L$ because
$\mathsf R_\delta V^{-1}=LA$. Hence $[z]\mapsto[Lz]$ is an explicit
isomorphism of the two cokernels. It is not by itself a map of tori or of
global threefolds.

\subsection{The original map's rational nonintersection}

Retain the exact form and orthogonal matrix
\[
 \Delta(q)=4(q_1q_2+q_1q_3+q_2q_3)-q_1^2-q_2^2-q_3^2,
 \quad
 F=\begin{pmatrix}
 1/\sqrt3&1/\sqrt3&1/\sqrt3\\
 1/\sqrt2&0&-1/\sqrt2\\
 1/\sqrt6&-2/\sqrt6&1/\sqrt6
 \end{pmatrix}.
\]
For $\det A=-\delta$, $\delta\in\mathbb Q_{>0}$, Pro's map is
\[
 \mathfrak K_\delta(A)=F^{-1}
 \frac{((a+c)/2,b,(a-c)/2)^{\mathsf T}}{\sqrt{3\delta}}.
\]
It satisfies $\Delta(\mathfrak K_\delta(A))=-1$ but sends no rational
$A$ to a rational triple. Indeed, if $q$ were rational, put
$S=q_1+q_2+q_3$, $U=q_1-q_3$, $V=q_1+q_3-2q_2$. The rational coordinates
$H=(a+c)/2$, $X=b$, $Y=(a-c)/2$ would satisfy
\[
 H=\sqrt\delta S,\quad X=\sqrt{3\delta/2}\,U,
 \quad Y=\sqrt{\delta/2}\,V.
\]
Two nonzero coordinates would make one of $\sqrt{3/2}$, $\sqrt2$,
$\sqrt3$ rational. If only $H$ is nonzero the determinant is positive.
If only $X$ or only $Y$ is nonzero, the determinant equation forces,
respectively, $U^2=2/3$ or $V^2=2$, impossible over $\mathbb Q$.
All coordinates zero also contradicts $\delta>0$. This proves the claim.

\subsection{A different rational bridge, with exact inverse and denominators}

The preceding nonintersection concerns the specified map, not every map
between the quadratic surfaces. The following construction uses projection
from the rational point $(1,0,0)$ on $\Delta=-1$; it does not replace $F$ or
the physical meaning, if any, assigned to its coordinates.

\begin{lemma}\label{lem:esv3-anisotropic}
The equation $\Delta(q)=0$ has no nonzero rational solution.
\end{lemma}
\begin{proof}
After clearing denominators use the integral $S,U,V$ just defined.
Expansion gives $2\Delta=2S^2-3U^2-V^2$. Modulo three,
$2S^2=V^2$ forces $3\mid S,V$. Substitution then forces $3\mid U$.
Repeating gives divisibility of $S,U,V$ by every power of three, hence all
vanish. The linear map from $q$ to $(S,U,V)$ has nonzero determinant,
so $q=0$.
\end{proof}

\begin{theorem}\label{thm:esv3-rational-bridge}
For fixed $\delta\in\mathbb Q_{>0}$ write
\[
 X_\delta=\{(a,b,c):ac-b^2=-\delta\},\quad
 Y=\{q:\Delta(q)=-1\},
\]
and define
\[
 D(a,b)=-1-a^2-b^2+4a+4b+4ab,\quad K(a,b)=1-2a-2b.
\]
On $U_\delta=\{(a,b,c)\in X_\delta:aDK\ne0\}$ the map
\[
 \mathcal B_\delta(a,b,c)=
 \left(1+\frac{2K}{D},\frac{2Ka}{D},\frac{2Kb}{D}\right)
 \label{eq:esv3-birational}
\]
is an isomorphism of open varieties over $\mathbb Q$ onto the open part of
$Y$ where $q_1\ne1$, $q_2\ne0$, and
$D(q_2/(q_1-1),q_3/(q_1-1))\ne0$. Its inverse is
\[
 a=\frac{q_2}{q_1-1},\qquad b=\frac{q_3}{q_1-1},\qquad
 c=\frac{b^2-\delta}{a}.
\]
Every integral matrix on $X_\delta$ with $a>0$ belongs to $U_\delta$.
For such a matrix, the exact least common denominator of its image is
\[
 d_{\rm image}=\frac{|D(a,b)|}{\gcd(|D(a,b)|,|2K(a,b)|)}.
\]
In particular its image is integral exactly when $D(a,b)\mid2K(a,b)$.
\end{theorem}
\begin{proof}
Set $q=(1,0,0)+\lambda(1,a,b)$. Expansion, with no change of scale in
$\Delta$, gives
\[
 \Delta(q)=-1-2K\lambda+D\lambda^2.
\]
The second intersection with $Y$ is therefore $\lambda=2K/D$, nonzero
on the stated domain. The displayed image and inverse follow. Conversely,
for a point of the stated target open set, put $\lambda=q_1-1$ and recover
$a,b$. Its equation gives $\lambda D=2K$, so $K\ne0$ and the constructions
are inverse regular maps on the indicated opens. Recovering $c$ retains the
fixed determinant $-\delta$. For integral $a,b$, $K$ is odd and cannot
vanish. Moreover $D=\Delta(1,a,b)\ne0$ by
Lemma~\ref{lem:esv3-anisotropic}. Finally, $q_1-1=2K/D$ has reduced
denominator $|D|/\gcd(|D|,|2K|)$, and $q_2,q_3$ are integer multiples
of this fraction. This proves both the denominator formula and its exactness.
\end{proof}

For the packet's fixed opening
$A_*=\left(\begin{smallmatrix}1&1\\1&0\end{smallmatrix}\right)$,
$\delta=1$, $D=9$, $K=-3$, and
\[
 \mathcal B_1(A_*)=(1/3,-2/3,-2/3),\qquad
 \Delta(1/3,-2/3,-2/3)=-1.
\]
This rational image coexists with the irrationality of
$\mathfrak K_1(A_*)$, because the two maps are different and both have now
been specified. The rational bridge does not automatically preserve the
integral positive chamber, the integral orbit count, or a dynamics. The
denominator condition above is the exact next arithmetic restriction, and
the inverse retains $\delta$ as a parameter rather than erasing it.

\subsection{Independent checks and reproducibility}

The separate verifier
\path{supporting_materials/workbench/evidence/koide_es_v3_audit_20260904/verify_independent.py}
enumerates the smaller denominator $N/R<y\le2N/R$, then reconstructs
$z=Ny/(Ry-N)$, independently of the packet's divisor enumerator. It agrees
with all 1430 oriented states in 665 coprime shells through $n=161$.
It checks the 819 diagonal certificates for 329 primes below 5000, the
symbolic rational-map identity and inverse, and 820 exact denominator cases.
These finite regressions supplement, rather than replace, the proofs above.
The original packet's verifier is replayed separately without modifying its
source or its original output. The archive SHA-256 is
\begin{quote}\small\ttfamily
ab5c21cf00268f9019c95ee24b565d05c4d9\allowbreak
337c51e4a576ded9e3cc3f0f235e
\end{quote}
The correct full hash is recorded byte-for-byte in the evidence manifest;
the line breaking here is typographical only.
